\documentclass[12pt,a4paper]{article}

\usepackage[T1]{fontenc}
\usepackage[utf8]{inputenc}
\usepackage[indonesian]{babel}
\usepackage[landscape,left=1.5cm,right=1.5cm,top=1.5cm,bottom=1.5cm]{geometry}
\usepackage{array}
\usepackage{tabularx}
\usepackage{times}

\pagestyle{plain}
\setlength{\parindent}{0pt}

\begin{document}

\begin{center}
    {\Large\bfseries Matriks Konsistensi Penelitian}\par
    \vspace{4pt}
    Nobel Nizam Fathirizki -- 123450117
\end{center}

\vspace{8pt}

\begin{table}[h]
    \centering
    \scriptsize
    \setlength{\tabcolsep}{3pt}
    \renewcommand{\arraystretch}{1.25}
    \begin{tabularx}{\linewidth}{|c|>{\raggedright\arraybackslash}X|>{\raggedright\arraybackslash}X|>{\raggedright\arraybackslash}X|>{\raggedright\arraybackslash}X|>{\raggedright\arraybackslash}X|}
        \hline
        \textbf{No.} & \textbf{Rumusan masalah} & \textbf{Tujuan penelitian} & \textbf{Data atau variabel} & \textbf{Metode analisis} & \textbf{Indikator dan luaran} \\
        \hline
        1 & Bagaimana perbedaan kinerja pipeline Medallion--Spark dan Pandas jika dievaluasi berdasarkan waktu pemrosesan dan throughput? & Mengevaluasi perbedaan kinerja kedua pipeline berdasarkan waktu pemrosesan dan throughput. & Rekaman kualitas air England; ukuran data; waktu mulai dan selesai; jumlah baris yang diproses. & Menjalankan proses \textit{batch} pada lingkungan lokal yang sama dan membandingkan hasil beberapa pengulangan. & Waktu pemrosesan ujung ke ujung, throughput (baris/detik), serta tabel perbandingan pipeline. \\
        \hline
        2 & Seberapa besar pengaruh penambahan kepadatan penduduk terhadap klasifikasi WQI jika dianalisis melalui perbandingan dengan parameter fisikokimia saja? & Menganalisis pengaruh \texttt{pop\_density} melalui perbandingan dengan model fisikokimia. & Delapan parameter fisikokimia; \texttt{pop\_density}; kelas WQI. & Studi ablasi: membandingkan model tanpa dan dengan \texttt{pop\_density} pada partisi data dan konfigurasi evaluasi yang sama. & Selisih akurasi dan F1-score makro, beserta analisis kontribusi fitur tanpa penafsiran kausal. \\
        \hline
        3 & Bagaimana perbedaan kinerja Random Forest dan XGBoost jika dibandingkan dalam klasifikasi WQI pada prosedur evaluasi yang sama? & Membandingkan kinerja Random Forest dan XGBoost pada prosedur evaluasi yang sama. & Tabel fitur lapisan \textit{Gold}; kelas WQI; keluaran prediksi kedua algoritma. & Melatih dan menguji kedua algoritma pada partisi data yang sama, kemudian membandingkan metrik klasifikasi per kelas dan secara keseluruhan. & Akurasi, F1-score makro, matriks konfusi, serta model dengan kinerja lebih baik pada data uji. \\
        \hline
    \end{tabularx}
\end{table}

\end{document}
